\chapter{Variables V01--V15 de priorización}
\label{cap:anexob}

Este anexo define las quince variables académicas utilizadas para estructurar la información del incidente y conectar la extracción de señales de Capa 1 con la priorización de Capa 2.

\section*{B.1. Finalidad del anexo}

Este anexo especifica las variables V01--V15 utilizadas como marco común entre datos,
etiquetas débiles, Capa 2, explicación y revisión humana. La matriz no debe interpretarse
como un formulario operativo oficial, sino como una representación académica de factores
relevantes para priorizar incidentes 112 en Castilla y León.

El alcance distingue entre variables implementadas de forma reproducible con el
Registro 112 CyL y variables previstas como ampliación futura. Esta distinción es clave para evitar
fuga de información y para no atribuir al modelo datos contextuales que todavía no se
han integrado.

\section*{B.2. Matriz de variables}

\begin{table}[ht]
\centering
\caption{Variables V01--V15 del prototipo}
\label{tab:anexob-variables}
\scriptsize
\begin{tabularx}{\textwidth}{|p{1.0cm}|p{3.0cm}|p{2.1cm}|X|p{2.1cm}|}
\hline
\textbf{ID} & \textbf{Variable} & \textbf{Tipo} & \textbf{Función en el sistema} & \textbf{Estado evaluado} \\ \hline
V01 & Riesgo vital inmediato & Binaria & Captura fallecido, inconsciente, herido grave, atrapado crítico o riesgo vital textual. & Implementada por señales \\ \hline
V02 & Víctimas confirmadas o probables & Ordinal & Aproxima magnitud humana inicial cuando el texto permite inferir afectados. & Parcial por texto \\ \hline
V03 & Tipo de daño principal & Categórica & Diferencia daño a personas, bienes, medioambiente o incidente menor. & Implementada indirecta \\ \hline
V04 & Tipo de incidente & Categórica & Mapea el incidente a la taxonomía del Anexo~\ref{cap:anexoa}. & Implementada \\ \hline
V05 & Población vulnerable & Binaria & Identifica menores, mayores, dependientes o colectivos especialmente expuestos. & Parcial por texto \\ \hline
V06 & Personas en riesgo estimadas & Ordinal & Representa exposición potencial más allá de víctimas confirmadas. & Parcial \\ \hline
V07 & Emplazamiento crítico & Categórica & Hospitales, residencias, centros educativos, infraestructuras o lugares sensibles. & Parcial por texto \\ \hline
V08 & Nivel de aviso AEMET & Ordinal & Contexto meteorológico oficial para elevar riesgo cuando sea pertinente. & Diferida \\ \hline
V09 & Fenómeno meteorológico activo & Categórica & Lluvia, nieve, viento, tormenta, temperatura u otro fenómeno coherente. & Diferida \\ \hline
V10 & Peligrosidad por inundación & Binaria & Cruce territorial con zonas inundables SNCZI. & Diferida \\ \hline
V11 & Proximidad Seveso & Binaria & Identifica instalaciones sujetas a RD 840/2015 o riesgo tecnológico. & Diferida \\ \hline
V12 & Tendencia del incidente & Ordinal & Estable, agravamiento, mejora o desconocida; orienta escalado. & Parcial por texto \\ \hline
V13 & Fiabilidad del informador & Ordinal & Modula confianza, no gravedad. & No disponible directa \\ \hline
V14 & Número de avisos simultáneos & Ordinal & Varias llamadas elevan confianza y posible impacto. & Implementada por señal \\ \hline
V15 & Accesibilidad al punto & Ordinal & Dificultad de acceso, entorno rural, montaña o rescate complejo. & Parcial por texto \\ \hline
\end{tabularx}
\end{table}

\section*{B.3. Variables implementadas en Capa 2}

El modelo RuleFit-lite seleccionado no consume directamente todas las variables V01--V15
como campos completos. Consume un conjunto compacto de señales predecisionales derivadas
del texto y de la categoría preliminar. Estas señales se obtienen mediante NLP
simbólico y auditable en Capa 1: detección semántica basada en reglas, tratamiento de
negaciones, normalización y confianza:

\begin{itemize}
    \item señales de fallecimiento, herido grave, atrapamiento e intoxicación;
    \item llamadas múltiples, incendio, accidente de tráfico, rescate y riesgo
    meteorológico o de inundación;
    \item riesgo vital expresado en el texto;
    \item variables categóricas normalizadas derivadas de la taxonomía.
\end{itemize}

Esta decisión mantiene reproducibilidad y evita introducir variables contextuales
incompletas. Las variables V08--V11 se conservan en el diseño porque son relevantes
para una ampliación futura, pero no se contabilizan como evidencia disponible en la
evaluación actual.

\section*{B.4. Columnas prohibidas por leakage y políticas de privacidad}

El sistema excluye rigurosamente del entrenamiento e inferencia cualquier variable o columna que refleje decisiones ex post o información administrativa y clínica posterior a la primera decisión de despacho. En particular, queda prohibido el uso de:

\begin{itemize}
    \item \texttt{MediosMov} y derivados de recursos reales movilizados.
    \item \texttt{PacientesAten} y registros de atención sanitaria o traslados.
    \item \texttt{IncidenteCerrado} y estado final de resolución del caso.
    \item \texttt{ultimaActualizacion} y marcas de tiempo de modificación del log central.
    \item Enlaces, identificadores de despacho y metadatos ex post.
\end{itemize}

La exclusión de estas columnas ex post posee una justificación metodológica crítica: si el modelo RuleFit tuviera acceso a variables como el número de ambulancias movilizadas o el estado de cierre del incidente, aprendería consecuencias operativas posteriores en lugar de señales tempranas de riesgo. Esto provocaría fuga de la variable objetivo (\textit{target leakage}) e invalidaría la utilidad del sistema en el momento de la primera decisión.

En alineación con los principios de minimización del Reglamento General de Protección de Datos (RGPD) \cite{RGPD} y de la Ley Orgánica 3/2018 (LOPDGDD) \cite{LOPDGDD}, el diseño contempla anonimización y tratamiento restringido de datos. Las coordenadas geográficas exactas del Registro 112 CyL pueden reducirse a precisión municipal para limitar el riesgo de reidentificación. Estas medidas son decisiones de diseño académico y no constituyen una certificación de conformidad regulatoria.

\section*{B.5. Ausencia de datos}

La ausencia de una variable no debe rellenarse con información posterior. Cuando el dato
no está disponible, el sistema debe:

\begin{itemize}
    \item emitir la recomendación con las señales presentes;
    \item reducir o matizar la confianza cuando falten variables relevantes;
    \item mostrar la ausencia en la explicación si afecta a la interpretación;
    \item registrar el caso para revisión si la ausencia coincide con un error crítico.
\end{itemize}

\section*{B.6. Limitaciones y riesgos}

La matriz V01--V15 es más amplia que la implementación evaluada. Esta decisión documenta la evolución futura sin atribuir al modelo información de AEMET, SNCZI o Seveso que todavía no utiliza.

Existe un riesgo crítico de leakage residual si en el futuro se introducen nuevas variables o integraciones del sistema de despacho sin verificar rigurosamente su disponibilidad temporal en la llamada de entrada. Cualquier campo que se incorpore al modelo debe someterse a una auditoría de línea de tiempo para asegurar que se genera en los primeros instantes de la llamada del ciudadano y no durante la posterior coordinación de la emergencia.
